\chapter{Future work}\label{chap:future-work}

Despite our efforts described in Chapter~\ref{chap:implementation}, we consider our implementation to be far from feature-complete. Therefore, we suggest several possible directions for future work.

\section{CSS support}

In \gls{svg}, elements can be styled using \gls{css}. \gls{css} properties have the same name as the corresponding \gls{svg} attributes and typically have the same syntax, but may occasionally offer more powerful functionality~\cite{svg11}. For this reason, we consider it valuable to implement proper support for parsing, manipulating, and serializing \gls{css} embedded in \gls{svg}.

\section{Optimization support}

Our implementation already provides a basic means of optimizing \gls{svg} documents using formatting options, reification, and rescaling of the coordinate system. However, these features need to be used manually, and the results strongly depend on the options used.

To address this limitation, we propose developing a clear and unified interface that would allow automatic optimization of \gls{svg} documents based on the application of various optimization techniques.

Along with this change, it is naturally desirable to extend our current set of optimization procedures with, for example, some of the following advanced techniques:
\begin{itemize}
    \item removing unused definitions,
    \item removing redundant attributes,
    \item removing invisible elements,
    \item simplifying and approximating path data,
    \item converting between basic shapes and their equivalent path representation (whichever is shorter),
    \item composing transformations in a transformation list into a single \texttt{matrix()} form.
\end{itemize}

We believe that the proposed changes would bring the library in line with other \gls{svg} optimization software, such as the popular \textsc{svgo}~\cite{svgo} project.
